\documentclass[11pt]{article}
\pagestyle{empty}

\setlength{\topmargin}{-.75in}
\setlength{\textheight}{9.2in}
\setlength{\oddsidemargin}{-.25in}
\setlength{\evensidemargin}{-.25in}
\setlength{\textwidth}{6.75in}
\setlength{\parskip}{4pt}

\usepackage{amsmath}
\usepackage{amsfonts}
\usepackage{amssymb}
\usepackage{booktabs}
\usepackage{graphicx}
\usepackage{mathtools}

\makeatletter
\renewcommand\section{\@startsection{section}{1}{\z@}%
  {-3.5ex \@plus -1ex \@minus -.2ex}%
  {2.3ex \@plus .2ex}%
  {\normalfont\Large}}
\renewcommand\subsection{\@startsection{subsection}{2}{\z@}%
  {-3.25ex\@plus -1ex \@minus -.2ex}%
  {1.5ex \@plus .2ex}%
  {\normalfont\large}}
\makeatother

\newcommand{\QOneGeneratedPath}{question1_output/question1_generated.tex}
\IfFileExists{\QOneGeneratedPath}{}{
  \renewcommand{\QOneGeneratedPath}{PSET_3/question1_output/question1_generated.tex}
}
\input{\QOneGeneratedPath}
\newcommand{\QThreeGeneratedPath}{question3_output/question3_generated.tex}
\newcommand{\QThreeSVARIRFPath}{question3_output/figures/q3_svar_iv_ebp_irf.pdf}
\newcommand{\QThreeFVDPath}{question3_output/figures/q3_svar_iv_ebp_fvd.pdf}
\newcommand{\QThreeIRFComparisonPath}{question3_output/figures/q3_ebp_irf_comparison.pdf}
\newcommand{\QThreeCIPath}{question3_output/figures/q3_optional_ci_panels.pdf}
\IfFileExists{\QThreeGeneratedPath}{}{
  \renewcommand{\QThreeGeneratedPath}{PSET_3/question3_output/question3_generated.tex}
  \renewcommand{\QThreeSVARIRFPath}{PSET_3/question3_output/figures/q3_svar_iv_ebp_irf.pdf}
  \renewcommand{\QThreeFVDPath}{PSET_3/question3_output/figures/q3_svar_iv_ebp_fvd.pdf}
  \renewcommand{\QThreeIRFComparisonPath}{PSET_3/question3_output/figures/q3_ebp_irf_comparison.pdf}
  \renewcommand{\QThreeCIPath}{PSET_3/question3_output/figures/q3_optional_ci_panels.pdf}
}
\input{\QThreeGeneratedPath}
\newcommand{\QFourGeneratedPath}{question4_output/question4_generated.tex}
\newcommand{\QFourFigurePath}{question4_output/figures/q4_sign_restricted_output_irf.pdf}
\IfFileExists{\QFourGeneratedPath}{}{
  \renewcommand{\QFourGeneratedPath}{PSET_3/question4_output/question4_generated.tex}
  \renewcommand{\QFourFigurePath}{PSET_3/question4_output/figures/q4_sign_restricted_output_irf.pdf}
}
\input{\QFourGeneratedPath}

\title{Advanced Time Series PSET 3}
\author{Dhruv Kohli}
\date{May 2026}

\begin{document}

\maketitle

\section*{Question 1}

The NKPC is
\[
\pi_t
=c+\gamma_f E_t(\pi_{t+1})+\gamma_b\pi_{t-1}+\lambda x_t+\nu_t,
\qquad
\gamma_b=1-\gamma_f,
\]
and inflation is assumed to satisfy
\[
\pi_t=\mu_\pi+\rho_\pi\pi_{t-1}+\rho_x x_{t-1}+\eta_t,
\qquad
(\nu_t,\eta_t)' \overset{i.i.d.}{\sim} (0,\Sigma).
\]
We measure inflation as
\[
\pi_t=100\Delta\log(\texttt{GDPDEF}_t),
\]
the quarter-over-quarter log change in the GDP deflator in percentage points.

\subsection*{(i)}
Since the information set at date \(t\) contains current and lagged \((\pi_t,x_t)\), rational
expectations and the inflation ADL imply
\[
E_t(\pi_{t+1})=\mu_\pi+\rho_\pi\pi_t+\rho_x x_t.
\]
Substituting this into the NKPC gives
\[
\pi_t
=c+\gamma_f(\mu_\pi+\rho_\pi\pi_t+\rho_x x_t)
+(1-\gamma_f)\pi_{t-1}+\lambda x_t+\nu_t.
\]
Let
\[
d(\theta)\equiv \lambda+\gamma_f\rho_x.
\]
When \(d(\theta)\neq 0\), rearranging the NKPC gives
\[
x_t
=
\frac{\mu_\pi(1-\gamma_f\rho_\pi-\gamma_f)-c}{d(\theta)}
+\frac{(1-\gamma_f\rho_\pi)\rho_\pi-(1-\gamma_f)}{d(\theta)}\pi_{t-1}
+\frac{(1-\gamma_f\rho_\pi)\rho_x}{d(\theta)}x_{t-1}
+\frac{(1-\gamma_f\rho_\pi)\eta_t-\nu_t}{d(\theta)}.
\]
Combining this equation with the ADL equation for \(\pi_t\), the vector
\[
y_t\equiv(\pi_t,x_t)'
\]
follows a bivariate VAR(1),
\[
y_t=a+A(\theta)y_{t-1}+u_t,
\]
where the four slope coefficients are
\[
A(\theta)
=
\begin{pmatrix}
\rho_\pi & \rho_x\\[4pt]
\dfrac{(1-\gamma_f\rho_\pi)\rho_\pi-(1-\gamma_f)}
{\lambda+\gamma_f\rho_x}
&
\dfrac{(1-\gamma_f\rho_\pi)\rho_x}
{\lambda+\gamma_f\rho_x}
\end{pmatrix}.
\]
The intercepts and innovations are
\[
a=
\begin{pmatrix}
\mu_\pi\\[3pt]
\dfrac{\mu_\pi(1-\gamma_f\rho_\pi-\gamma_f)-c}{\lambda+\gamma_f\rho_x}
\end{pmatrix},
\qquad
u_t=
\begin{pmatrix}
\eta_t\\[3pt]
\dfrac{(1-\gamma_f\rho_\pi)\eta_t-\nu_t}{\lambda+\gamma_f\rho_x}
\end{pmatrix}.
\]
Because \((\nu_t,\eta_t)'\) is i.i.d., \(u_t\) is also an i.i.d. reduced-form innovation.

\subsection*{(ii)}
We estimate the unrestricted reduced-form VAR(1)
\[
y_t=a+A y_{t-1}+u_t
\]
by least squares using the same source data and sample as PSET 2, Question 6:
\texttt{GDPDEF} and quarterly-average \texttt{UNRATE}, restricted to the largest
overlapping sample before 2020 and using the percentage-point inflation scaling stated above.
The raw sample is \QOneRawStart--\QOneRawEnd\ (\QOneRawN\ quarters), and the VAR
regression sample is \QOneVARStart--\QOneVAREnd\ (\QOneVARN\ observations).

The estimated intercepts are
\[
\hat a_\pi=\QOneInterceptPi,
\qquad
\hat a_x=\QOneInterceptX.
\]
The estimated slope matrix is
\[
\widehat A
=
\begin{pmatrix}
\QOneAOneOne & \QOneAOneTwo\\
\QOneATwoOne & \QOneATwoTwo
\end{pmatrix},
\]
with rows corresponding to the \(\pi_t\) and \(x_t\) equations and columns corresponding
to \(\pi_{t-1}\) and \(x_{t-1}\). The largest absolute eigenvalue of \(\widehat A\) is
\(\QOneMaxEigen\).

The minimum-distance estimator uses the restriction
\[
\widehat\theta
=
\arg\min_{\theta}
\{\operatorname{vec}(\widehat A)-\operatorname{vec}(A(\theta))\}'
W
\{\operatorname{vec}(\widehat A)-\operatorname{vec}(A(\theta))\}.
\]
We use \(W=I_4\). Since the model gives four reduced-form slope coefficients for the four
parameters \(\theta=(\gamma_f,\lambda,\rho_\pi,\rho_x)'\), the estimate is exactly
identified. In particular,
\[
\widehat\rho_\pi=\widehat A_{11},
\qquad
\widehat\rho_x=\widehat A_{12},
\]
and, using the last two VAR coefficients,
\[
\widehat\gamma_f
=
\frac{\widehat A_{22}(1-\widehat\rho_\pi)+\widehat A_{21}\widehat\rho_x}
{\widehat A_{21}\widehat\rho_\pi\widehat\rho_x
+\widehat A_{22}(1-\widehat\rho_\pi^2)},
\]
\[
\widehat\lambda
=
\frac{(1-\widehat\gamma_f\widehat\rho_\pi)\widehat\rho_x}{\widehat A_{22}}
-\widehat\gamma_f\widehat\rho_x.
\]
The resulting estimates are:
\[
\begin{array}{lccc}
\toprule
\text{Parameter} & \text{Estimate} & \text{95\% bootstrap lower} & \text{95\% bootstrap upper}\\
\midrule
\gamma_f & \QOneGammaF & \QOneGammaFCILower & \QOneGammaFCIUpper\\
\lambda & \QOneLambda & \QOneLambdaCILower & \QOneLambdaCIUpper\\
\rho_\pi & \QOneRhoPi & \QOneRhoPiCILower & \QOneRhoPiCIUpper\\
\rho_x & \QOneRhoX & \QOneRhoXCILower & \QOneRhoXCIUpper\\
\bottomrule
\end{array}
\]

\subsection*{(iii)}
We compute the confidence intervals using a recursive residual bootstrap for the fitted
VAR(1). The procedure is:

\begin{enumerate}
\item Estimate the VAR(1) by OLS and collect the centered residual vectors
\(\hat u_t-\bar{\hat u}\).
\item For each bootstrap draw \(b=1,\ldots,\QOneBootstrapReps\), sample residual vectors
with replacement from the centered empirical residual distribution, preserving the
cross-equation dependence between the \(\pi_t\) and \(x_t\) shocks.
\item Starting from the observed initial value \(y_0^*=y_0\), recursively generate
\[
y_t^*=\hat a+\widehat A y_{t-1}^*+u_t^*.
\]
\item Re-estimate the VAR(1) on the bootstrap sample and recompute
\(\widehat\theta^*=(\widehat\gamma_f^*,\widehat\lambda^*,
\widehat\rho_\pi^*,\widehat\rho_x^*)'\) using the same minimum-distance inversion.
\item Report the 2.5th and 97.5th percentiles of the bootstrap estimates.
\end{enumerate}

We use \(\QOneBootstrapReps\) bootstrap repetitions with random seed \(\QOneBootstrapSeed\);
all \(\QOneBootstrapUsed\) repetitions produced nonsingular estimates. This bootstrap relies
on the VAR(1) restriction derived above, stationarity of the fitted VAR, and the i.i.d.
innovation assumption. The i.i.d. assumption is exactly the extra structure imposed in the
question, since the reduced-form VAR innovations are linear transformations of
\((\nu_t,\eta_t)'\).

\section*{Question 2}

We observe
\[
y_t=\mu+z_t,
\]
where \(\mu\) is unknown and \(\{z_t\}\) is mean zero, strictly stationary, and weakly
dependent. The sample mean is
\[
\widehat\mu=\frac{1}{T}\sum_{t=1}^T y_t.
\]
The long-run variance needed for inference on \(\mu\) is
\[
\Omega=\operatorname{LRV}(z_t)=\sum_{\ell=-\infty}^{\infty}\Gamma_z(\ell).
\]
Since \(y_t\) differs from \(z_t\) only by a constant, the population autocovariances of
\(\{y_t\}\) and \(\{z_t\}\) are the same. The kernel HAC estimator is
\[
\widehat\Omega
=
\sum_{\ell=-(T-1)}^{T-1}
k\!\left(\frac{\ell}{S_T}\right)\widehat\Gamma(\ell),
\]
where, for \(\ell\geq 0\),
\[
\widehat\Gamma(\ell)
=
\frac{1}{T}\sum_{t=\ell+1}^T
(y_t-\widehat\mu)(y_{t-\ell}-\widehat\mu),
\qquad
\widehat\Gamma(-\ell)=\widehat\Gamma(\ell)
\]
in this scalar problem.

Under the Andrews assumptions, the one-dimensional MSE-optimal
bandwidth is
\[
S_T^*
=
\left[
\frac{q^* k_{q^*}^2}{\int_{-\infty}^{\infty} k(x)^2\,dx}
\alpha(q^*)T
\right]^{1/(2q^*+1)},
\]
where
\[
\alpha(q)=
\left(\frac{s_z^{(q)}(0)}{\Omega}\right)^2,
\qquad
s_z^{(q)}(0)=
\sum_{\ell=-\infty}^{\infty}|\ell|^q\Gamma_z(\ell).
\]

\subsection*{(i)}
Under the stationary AR(1) working model
\[
z_t=\rho z_{t-1}+\varepsilon_t,
\qquad
\varepsilon_t\sim WN(0,\sigma^2),
\qquad
|\rho|<1.
\]
Its autocovariance function is
\[
\Gamma_z(\ell)=\frac{\sigma^2}{1-\rho^2}\rho^{|\ell|}.
\]
Therefore
\[
\Omega
=
\frac{\sigma^2}{1-\rho^2}
\left(1+2\sum_{\ell=1}^{\infty}\rho^\ell\right)
=
\frac{\sigma^2}{1-\rho^2}
\left(\frac{1+\rho}{1-\rho}\right)
=
\frac{\sigma^2}{(1-\rho)^2}.
\]
For any \(q>0\), the zero-lag term in \(s_z^{(q)}(0)\) is zero, and since
\(|\rho|<1\) the following series is absolutely convergent:
\[
s_z^{(q)}(0)
=
2\sum_{\ell=1}^{\infty}\ell^q\Gamma_z(\ell)
=
\frac{2\sigma^2}{1-\rho^2}
\sum_{\ell=1}^{\infty}\ell^q\rho^\ell.
\]
Thus
\[
\frac{s_z^{(q)}(0)}{\Omega}
=
\frac{2(1-\rho)}{1+\rho}
\sum_{\ell=1}^{\infty}\ell^q\rho^\ell,
\]
and
\[
\alpha(q)
=
\left[
\frac{2(1-\rho)}{1+\rho}
\sum_{\ell=1}^{\infty}\ell^q\rho^\ell
\right]^2.
\]
Substituting this into Andrews's formula gives
\[
S_T^*
=
\left[
\frac{q^*k_{q^*}^2T}{\int_{-\infty}^{\infty} k(x)^2\,dx}
\left\{
\frac{2(1-\rho)}{1+\rho}
\sum_{\ell=1}^{\infty}\ell^{q^*}\rho^\ell
\right\}^2
\right]^{1/(2q^*+1)}.
\]
This depends only on \(T\), \(\rho\), and the chosen kernel. The innovation variance
\(\sigma^2\) cancels, as it should: the plug-in bandwidth is determined by persistence,
not by the scale of the series.

\subsection*{(ii)}
The Bartlett kernel is
\[
k(x)=\max\{1-|x|,0\}.
\]
Near zero, \(k(x)=1-|x|\), so
\[
k_1=\lim_{x\to 0}\frac{1-k(x)}{|x|}=1.
\]
For \(q>1\),
\[
\frac{1-k(x)}{|x|^q}=|x|^{1-q}\to\infty,
\]
so the characteristic exponent is \(q^*=1\). Also
\[
\int_{-\infty}^{\infty} k(x)^2\,dx
=
\int_{-1}^{1}(1-|x|)^2\,dx
=
2\int_0^1(1-x)^2\,dx
=
\frac{2}{3}.
\]
For \(q=1\),
\[
\sum_{\ell=1}^{\infty}\ell\rho^\ell
=
\frac{\rho}{(1-\rho)^2},
\]
so
\[
\frac{s_z^{(1)}(0)}{\Omega}
=
\frac{2(1-\rho)}{1+\rho}
\frac{\rho}{(1-\rho)^2}
=
\frac{2\rho}{1-\rho^2},
\]
and
\[
\alpha(1)=\frac{4\rho^2}{(1-\rho^2)^2}.
\]
Substituting \(q^*=1\), \(k_1=1\), and \(\int k^2=2/3\) into the optimal bandwidth
formula gives
\[
S_{T,\mathrm{Bartlett}}^*
=
\left[
\frac{1}{2/3}
\frac{4\rho^2}{(1-\rho^2)^2}T
\right]^{1/3}
=
\left[
\frac{6\rho^2}{(1-\rho^2)^2}T
\right]^{1/3}.
\]
Equivalently,
\[
S_{T,\mathrm{Bartlett}}^*
=
\left(\frac{6\rho^2}{(1-\rho^2)^2}\right)^{1/3}T^{1/3}.
\]
For example, if \(\rho=0.25\), this gives approximately \(0.75T^{1/3}\), matching the
standard rule of thumb.

For the quadratic spectral kernel,
\[
k(x)
=
\frac{25}{12\pi^2x^2}
\left[
\frac{\sin(6\pi x/5)}{6\pi x/5}
-\cos(6\pi x/5)
\right],
\qquad x\neq 0,
\]
with \(k(0)=1\). The problem states that
\[
\int_{-\infty}^{\infty}k(x)^2\,dx=1,
\qquad
q^*=2,
\qquad
k_2=1.421223.
\]
To verify the local expansion, set \(a=6\pi/5\). Since
\[
\frac{\sin(ax)}{ax}=1-\frac{a^2x^2}{6}+\frac{a^4x^4}{120}+O(x^6),
\qquad
\cos(ax)=1-\frac{a^2x^2}{2}+\frac{a^4x^4}{24}+O(x^6),
\]
we have
\[
\frac{\sin(ax)}{ax}-\cos(ax)
=
\frac{a^2x^2}{3}-\frac{a^4x^4}{30}+O(x^6).
\]
Multiplying by \(25/(12\pi^2x^2)\) gives
\[
k(x)=1-\frac{18\pi^2}{125}x^2+O(x^4),
\]
so
\[
k_2=\frac{18\pi^2}{125}=1.421223\ldots,
\qquad q^*=2.
\]

For \(q=2\),
\[
\sum_{\ell=1}^{\infty}\ell^2\rho^\ell
=
\frac{\rho(1+\rho)}{(1-\rho)^3},
\]
and therefore
\[
\frac{s_z^{(2)}(0)}{\Omega}
=
\frac{2(1-\rho)}{1+\rho}
\frac{\rho(1+\rho)}{(1-\rho)^3}
=
\frac{2\rho}{(1-\rho)^2}.
\]
Thus
\[
\alpha(2)=\frac{4\rho^2}{(1-\rho)^4}.
\]
Substituting \(q^*=2\), \(k_2=1.421223\), and \(\int k^2=1\),
\[
S_{T,\mathrm{QS}}^*
=
\left[
2(1.421223)^2
\frac{4\rho^2}{(1-\rho)^4}T
\right]^{1/5}
=
\left[
\frac{8(1.421223)^2\rho^2}{(1-\rho)^4}T
\right]^{1/5}.
\]
Equivalently,
\[
S_{T,\mathrm{QS}}^*
=
1.3221
\left[
\frac{4\rho^2}{(1-\rho)^4}T
\right]^{1/5},
\qquad
1.3221=\{2(1.421223)^2\}^{1/5}.
\]

The difference in rates is coming from the kernel smoothness at the origin:
Bartlett has \(q^*=1\), so its bandwidth is proportional to \(T^{1/3}\), while QS has
\(q^*=2\), so its bandwidth is proportional to \(T^{1/5}\). The denominators also line
up with the AR(1) autocovariances: Bartlett depends on \(1-\rho^2\) because
\(s_z^{(1)}(0)/\Omega=2\rho/(1-\rho^2)\), whereas QS depends on \((1-\rho)^4\)
because \(s_z^{(2)}(0)/\Omega=2\rho/(1-\rho)^2\). If \(\rho=0\), the leading bias
term in the Andrews expansion is zero, so the plug-in formulas return \(S_T^*=0\).
In practice we would impose an integer bandwidth rule, but this white-noise edge case
is a feature of the plug-in MSE calculation rather than an algebraic problem.

\section*{Question 3}

We collect the four macro-financial variables in
\[
y_t=(\texttt{ipgr}_t,\texttt{infl}_t,\texttt{gs1}_t,\texttt{ebp}_t)'.
\]
The external instrument is \(z_t=\texttt{ff4\_tc}_t\),
the high-frequency change in three-month-ahead federal funds futures prices around
FOMC announcements. This is the same type of high-frequency policy surprise emphasized
by Gertler and Karadi (2015): it uses variation in futures
prices in a narrow policy-announcement window as an external source of monetary policy
variation, avoiding the simultaneity problem that arises when financial variables enter
the VAR. Following the question, we include an intercept and \(\QThreeLags\) lags in
every VAR and regression. The one-year Treasury rate, \(\texttt{gs1}\), is the policy
indicator, and all responses below are normalized so that the impact response of
\(\texttt{gs1}\) equals \(1\) (100 basis points).

\subsection*{(i)}
We estimate the reduced-form VAR
\[
y_t=c+\sum_{j=1}^{p}A_j y_{t-j}+\eta_t,
\qquad p=\QThreeLags,
\]
by least squares. Let
\[
\Sigma=E(\eta_t\eta_t'),
\qquad
\gamma=E(\eta_t z_t).
\]
Under the SVAR-IV assumptions,
\[
\eta_t=H\varepsilon_t,
\qquad
E(z_t\varepsilon_{1t})\neq 0,
\qquad
E(z_t\varepsilon_{jt})=0\quad(j\geq 2),
\]
with \(E(\varepsilon_t\varepsilon_t')=I\). Therefore
\(\gamma=h_1 E(z_t\varepsilon_{1t})\), where \(h_1\) is the first column of \(H\), so
\(\gamma\) is proportional to the impact vector of the monetary policy shock. With the
unit-variance structural-shock
normalization,
\[
h_1
=
\frac{\gamma}{(\gamma'\Sigma^{-1}\gamma)^{1/2}},
\]
up to sign. We choose the sign so the \(\texttt{gs1}\) impact response is positive. The
estimated first-stage regression of the \(\texttt{gs1}\) VAR residual on \(\texttt{ff4\_tc}\)
has \(F=\QThreeSVARFirstStageF\) and \(R^2=\QThreeSVARFirstStageRTwo\%\), so the
instrument is relevant in this sample.

Let \(\Psi_h\) denote the VAR moving-average coefficient matrices, with
\[
\Psi_0=I,
\qquad
\Psi_h=\sum_{j=1}^{\min(p,h)}A_j\Psi_{h-j}.
\]
The normalized impulse response is
\[
\operatorname{IRF}_h
=
\Psi_h\frac{h_1}{e_{\texttt{gs1}}'h_1},
\qquad h=0,1,\ldots,\QThreeHorizon.
\]
The resulting EBP response is shown in Figure~\ref{fig:q3-svar-irf}. The impact
response is \(\QThreeSVAREBPHZero\) percentage points. It remains positive for the whole
36-month window, falling to \(\QThreeSVAREBPHTwelve\) after 12 months and
\(\QThreeSVAREBPHThirtySix\) after 36 months. The sign and magnitude are consistent with the
main Gertler-Karadi mechanism: a contractionary policy surprise raises credit costs,
and the excess bond premium moves by a large amount relative to the one-year rate
normalization.

\begin{figure}[htbp]
\centering
\includegraphics[width=0.78\textwidth]{\QThreeSVARIRFPath}
\caption{SVAR-IV impulse response of the excess bond premium to a monetary policy
shock that raises \(\texttt{gs1}\) by 100 basis points on impact.}
\label{fig:q3-svar-irf}
\end{figure}

\subsection*{(ii)}
For the EBP forecast variance decomposition, use the unit-variance impact vector \(h_1\),
not the 100-basis-point rescaled vector. For horizon \(\ell\geq 1\),
\[
\operatorname{FVD}_{\texttt{ebp},1,\ell}
=
\frac{\sum_{h=0}^{\ell-1}
\left(e_{\texttt{ebp}}'\Psi_h h_1\right)^2}
{\sum_{h=0}^{\ell-1}
e_{\texttt{ebp}}'\Psi_h\Sigma\Psi_h'e_{\texttt{ebp}}}.
\]
The numerator is the forecast-error variance contribution of the monetary shock, and
the denominator is the total reduced-form VAR forecast-error variance. The estimated
FVD is plotted in Figure~\ref{fig:q3-fvd}. Monetary policy shocks explain about
\(\QThreeFVDHOne\%\) of the one-month EBP forecast-error variance,
\(\QThreeFVDHTwelve\%\) at 12 months, and \(\QThreeFVDHThirtySix\%\) at 36 months.

\begin{figure}[htbp]
\centering
\includegraphics[width=0.78\textwidth]{\QThreeFVDPath}
\caption{SVAR-IV forecast variance decomposition of the excess bond premium with
respect to the monetary policy shock.}
\label{fig:q3-fvd}
\end{figure}

\subsection*{(iii)}
We next drop the invertibility assumption. Following the standard LP-IV argument, form
\[
w_t=(z_t,y_t')'
=
(\texttt{ff4\_tc}_t,\texttt{ipgr}_t,\texttt{infl}_t,\texttt{gs1}_t,\texttt{ebp}_t)'.
\]
We estimate a VAR(\(\QThreeLags\)) for \(w_t\), order \(z_t\) first, and use a recursive
Cholesky shock to the first variable. Since \(z_t\) is only a noisy proxy for the monetary
shock, the absolute scale of this recursive shock is not structural. However, the
attenuation from measurement error is common across horizons and variables, so the
relative IRF is valid after normalizing by the impact response of \(\texttt{gs1}\). Thus we
again divide the responses by the impact response of \(\texttt{gs1}\).

The recursive-VAR LP-IV estimate is close to the SVAR-IV estimate on impact:
\(\QThreeLPVAREBPHZero\) percentage points for EBP. It then decays somewhat faster,
reaching \(\QThreeLPVAREBPHTwelve\) at 12 months and \(\QThreeLPVAREBPHThirtySix\) at 36 months.
The comparison is shown in Figure~\ref{fig:q3-irf-comparison}.

\subsection*{(iv)}
The direct 2SLS implementation estimates, for each horizon \(h=0,\ldots,\QThreeHorizon\),
\[
\texttt{ebp}_{t+h}
=
\alpha_h+\beta_h\texttt{gs1}_t
+\delta_h'X_{t-1}+u_{t+h|t},
\]
using \(z_t=\texttt{ff4\_tc}_t\) as an instrument for \(\texttt{gs1}_t\). The control vector
\(X_{t-1}\) contains an intercept, \(\QThreeLags\) lags of \(y_t\), and \(\QThreeLags\)
lags of \(z_t\). With a single endogenous regressor and a single IV, the coefficient can
be written after residualizing on controls as
\[
\widehat\beta_h
=
\frac{\tilde z'\tilde y_h}{\tilde z'\widetilde{\texttt{gs1}}},
\]
where tildes denote residuals after projection on \(X_{t-1}\). The first-stage \(F\)-statistic
at \(h=0\) is \(\QThreeLPFirstStageF\), with partial \(R^2=\QThreeLPPartialRTwo\%\).

The 2SLS LP-IV estimates are more jagged than the VAR-based estimates, which is expected
because each horizon is estimated separately and no VAR propagation restrictions are
imposed. The impact response is mechanically the same as the recursive-VAR LP-IV impact
response, \(\QThreeLPTwoSLSEBPHZero\). At longer horizons the direct estimates are larger:
\(\QThreeLPTwoSLSEBPHTwelve\) at 12 months and \(\QThreeLPTwoSLSEBPHThirtySix\) at 36 months.

\begin{figure}[htbp]
\centering
\includegraphics[width=0.78\textwidth]{\QThreeIRFComparisonPath}
\caption{EBP impulse responses from SVAR-IV, recursive-VAR LP-IV, and direct 2SLS
LP-IV. All responses are normalized to a 100 basis point impact response of
\(\texttt{gs1}\).}
\label{fig:q3-irf-comparison}
\end{figure}

\begin{center}
\begin{tabular}{lccc}
\toprule
Horizon & SVAR-IV & LP-IV recursive VAR & LP-IV 2SLS\\
\midrule
0 & \QThreeSVAREBPHZero & \QThreeLPVAREBPHZero & \QThreeLPTwoSLSEBPHZero\\
6 & \QThreeSVAREBPHSix & \QThreeLPVAREBPHSix & \QThreeLPTwoSLSEBPHSix\\
12 & \QThreeSVAREBPHTwelve & \QThreeLPVAREBPHTwelve & \QThreeLPTwoSLSEBPHTwelve\\
24 & \QThreeSVAREBPHTwentyFour & \QThreeLPVAREBPHTwentyFour & \QThreeLPTwoSLSEBPHTwentyFour\\
36 & \QThreeSVAREBPHThirtySix & \QThreeLPVAREBPHThirtySix & \QThreeLPTwoSLSEBPHThirtySix\\
\bottomrule
\end{tabular}
\end{center}

\subsection*{(v)}
We know the key testable implication of invertibility with an external IV:
if the monetary policy shock is invertible with respect to \(\{y_t\}\), then the proxy
\(z_t\) should not Granger-cause \(y_t\) conditional on lags of \(y_t\). Intuitively,
under invertibility the useful shock component of \(z_t\) is already contained in the
history of \(y_t\), while the measurement-error component of \(z_t\) is not useful for
forecasting.

Using the VAR specification in part (iii), we test the joint null that lags 1 through
\(\QThreeLags\) of \(\texttt{ff4\_tc}\) have zero coefficients in the four \(y_t\) equations.
This imposes \(4\times \QThreeLags=\QThreeInvertDF\) restrictions. The Gaussian
VAR likelihood-ratio statistic is
\[
LR
=
T\left(\log|\widehat\Sigma_R|-\log|\widehat\Sigma_U|\right),
\]
where \(\widehat\Sigma_U\) is the residual covariance matrix from the unrestricted
\(y_t\) equations with lags of \((z_t,y_t')'\), and \(\widehat\Sigma_R\) is from the
restricted equations with only lags of \(y_t\). The result is
\[
LR=\QThreeInvertLR,
\qquad
df=\QThreeInvertDF,
\qquad
p\text{-value}=\QThreeInvertPValue.
\]
Thus we do not reject the no-Granger-causality implication of invertibility. This does not
prove invertibility, but it means this particular testable implication is not contradicted
by this VAR specification.

\subsection*{(vi)}
We compute confidence intervals by a recursive residual bootstrap for the joint VAR in
\[
w_t=(z_t,y_t')'.
\]
This is the reduced-form object used in the recursive-VAR LP-IV estimator, and it also
gives the time-series model for \(z_t\) that is needed when bootstrapping SVAR-IV with
an external instrument. After estimating the VAR(\(\QThreeLags\)) for \(w_t\), we center
the joint residuals, resample the residual vectors with replacement, hold the first
\(\QThreeLags\) observations fixed, and recursively generate
\[
w_t^*
=\widehat c+\sum_{j=1}^{\QThreeLags}\widehat A_j w_{t-j}^*+u_t^*.
\]
For each artificial sample, we re-estimate the SVAR-IV IRF, the SVAR-IV FVD, the
recursive-VAR LP-IV IRF, and the direct 2SLS LP-IV IRF exactly as above. This means the
bootstrap includes uncertainty about the VAR propagation, the reduced-form covariance
matrix, the proxy covariance
\[
\widehat{\gamma}=\frac{1}{T}\sum_t\widehat{\eta}_t z_t,
\]
and the first-stage relationship used in the 2SLS local projections.

For a scalar object \(\theta_h\), we report the equal-tailed percentile interval
\[
\left[
q_{0.05,h}^*,
q_{0.95,h}^*
\right],
\]
where \(q_{\alpha,h}^*\) is the \(\alpha\)-quantile of the bootstrap estimates
\(\widehat\theta_h^*\). We use \QThreeBootstrapReps\ bootstrap repetitions with random
seed \(\QThreeBootstrapSeed\), and all \QThreeBootstrapUsed\ repetitions produced finite
estimates. The validity of this procedure relies on the maintained finite-order VAR
approximation for \(w_t\), stable dynamics, and i.i.d. joint VAR innovations. It is the
standard strong-instrument bootstrap approximation rather than a weak-IV-robust
procedure; the first-stage statistics above are why this is a reasonable approximation
here.

The selected 90\% confidence intervals for the EBP impulse responses are:
\begin{center}
\begin{tabular}{lccc}
\toprule
Estimator & \(h=0\) & \(h=12\) & \(h=36\)\\
\midrule
SVAR-IV &
\(\QThreeSVAREBPHZero\,[\QThreeCISVAREBPLowerZero,\QThreeCISVAREBPUpperZero]\) &
\(\QThreeSVAREBPHTwelve\,[\QThreeCISVAREBPLowerTwelve,\QThreeCISVAREBPUpperTwelve]\) &
\(\QThreeSVAREBPHThirtySix\,[\QThreeCISVAREBPLowerThirtySix,\QThreeCISVAREBPUpperThirtySix]\)\\
LP-IV recursive VAR &
\(\QThreeLPVAREBPHZero\,[\QThreeCILPVAREBPLowerZero,\QThreeCILPVAREBPUpperZero]\) &
\(\QThreeLPVAREBPHTwelve\,[\QThreeCILPVAREBPLowerTwelve,\QThreeCILPVAREBPUpperTwelve]\) &
\(\QThreeLPVAREBPHThirtySix\,[\QThreeCILPVAREBPLowerThirtySix,\QThreeCILPVAREBPUpperThirtySix]\)\\
LP-IV 2SLS &
\(\QThreeLPTwoSLSEBPHZero\,[\QThreeCILPTwoSLSEBPLowerZero,\QThreeCILPTwoSLSEBPUpperZero]\) &
\(\QThreeLPTwoSLSEBPHTwelve\,[\QThreeCILPTwoSLSEBPLowerTwelve,\QThreeCILPTwoSLSEBPUpperTwelve]\) &
\(\QThreeLPTwoSLSEBPHThirtySix\,[\QThreeCILPTwoSLSEBPLowerThirtySix,\QThreeCILPTwoSLSEBPUpperThirtySix]\)\\
\bottomrule
\end{tabular}
\end{center}
For the EBP forecast variance decomposition, the selected 90\% confidence intervals,
reported in percent, are:
\begin{center}
\begin{tabular}{lccc}
\toprule
 & \(\ell=1\) & \(\ell=12\) & \(\ell=36\)\\
\midrule
SVAR-IV FVD &
\(\QThreeFVDHOne\,[\QThreeCIFVDLowerOne,\QThreeCIFVDUpperOne]\) &
\(\QThreeFVDHTwelve\,[\QThreeCIFVDLowerTwelve,\QThreeCIFVDUpperTwelve]\) &
\(\QThreeFVDHThirtySix\,[\QThreeCIFVDLowerThirtySix,\QThreeCIFVDUpperThirtySix]\)\\
\bottomrule
\end{tabular}
\end{center}
Figure~\ref{fig:q3-ci-panels} plots the point estimates and intervals across all horizons.
The direct 2SLS intervals are much wider than the two VAR-based intervals, especially at
medium and long horizons. That is the empirical cost of estimating each horizon separately
without imposing the VAR propagation restrictions. The SVAR-IV and recursive-VAR LP-IV
bands are tighter, but they are also more model-dependent because their horizon profiles
come from the estimated VAR law of motion.

\begin{figure}[htbp]
\centering
\includegraphics[width=0.90\textwidth]{\QThreeCIPath}
\caption{Point estimates and 90\% recursive-residual bootstrap confidence intervals for
the SVAR-IV IRF, SVAR-IV FVD, recursive-VAR LP-IV IRF, and direct 2SLS LP-IV IRF.}
\label{fig:q3-ci-panels}
\end{figure}

\clearpage
\section*{Question 4}

We use the FRED series \texttt{INDPRO}, \texttt{PCEPI}, and \texttt{GS1}. Output growth
and inflation are monthly log growth rates in annualized percentage points,
\[
g_t=1200\Delta\log(\texttt{INDPRO}_t),
\qquad
\pi_t=1200\Delta\log(\texttt{PCEPI}_t),
\]
and the real interest rate is
\[
r_t=\texttt{GS1}_t-\pi_t.
\]
After constructing these variables, we keep observations dated \QFourSampleStart--\QFourSampleEnd,
so the VAR sample after \(\QFourLags\) lags is \QFourVARStart--\QFourVAREnd\
(\(\QFourVARN\) observations). Let
\[
y_t=(g_t,\pi_t,r_t)'.
\]

We estimate the reduced-form VAR(\(\QFourLags\))
\[
y_t=c+\sum_{j=1}^{\QFourLags} A_j y_{t-j}+\eta_t,
\qquad
E(\eta_t\eta_t')=\Sigma,
\]
by least squares. The largest absolute eigenvalue of the companion matrix is
\(\QFourMaxEigen\), so the estimated propagation is stable. Using the standard state-space
notation, define the companion state
\[
s_t=(y_t',y_{t-1}',\ldots,y_{t-\QFourLags+1}')',
\]
and write deviations from the deterministic intercept as
\[
s_t=F s_{t-1}+G\eta_t,\qquad y_t=J s_t.
\]
Here \(F\) contains \(A_1,\ldots,A_{\QFourLags}\) in its first block row and identity
matrices below the diagonal, \(G=(I_3,0,\ldots,0)'\), and \(J=(I_3,0,\ldots,0)\).
Thus the reduced-form moving-average matrices are
\[
\Psi_\ell=J F^\ell G,\qquad \ell=0,1,\ldots,\QFourHorizon.
\]
This is the standard linear state-space state-transition logic: an impulse to
the innovation changes the state through \(G\), propagates through \(F^\ell\), and is read
off by \(J\).

Let \(C\) be the lower Cholesky factor of \(\widehat\Sigma\), so \(CC'=\widehat\Sigma\).
Every exactly identified orthogonalization can be written as
\[
\eta_t=CQ\varepsilon_t,\qquad QQ'=I_3,\qquad E(\varepsilon_t\varepsilon_t')=I_3.
\]
For a candidate first column \(q=Q e_1\), the impulse response to the first
unit-variance structural shock is
\[
\theta_\ell(q)=e_g'\Psi_\ell Cq,\qquad \ell=0,1,\ldots,36,
\]
where \(e_g\) selects output growth. The contractionary monetary shock is restricted to
raise the real rate and lower inflation at horizons \(0,1,2\):
\[
e_r'\Psi_h Cq\geq 0,
\qquad
e_\pi'\Psi_h Cq\leq 0,
\qquad h=0,1,2.
\]
There is no sign restriction on output growth and no restriction on the other two shocks.
Because Haar measure is invariant to relabeling columns, restricting the first column is
without loss of generality for the one-shock object. We draw \QFourRotationDraws\
orthogonal matrices from Haar measure using QR decompositions of Gaussian matrices,
normalizing the diagonal of \(R\) to be positive as required in the assignment. Of these,
\QFourAcceptedDraws\ draws satisfy the restrictions, for an acceptance rate of
\QFourAcceptanceRate\%.

\subsection*{(i)}
The estimated identified-set bounds for \(\theta_\ell\) are the pointwise minimum and
maximum values of \(\theta_\ell(q)\) across the retained rotations. Selected horizons are:
\[
\begin{array}{lrrrr}
\toprule
\text{Horizon} & \text{Lower bound} & \text{Upper bound} & \text{Midpoint} & \text{Posterior median}\\
\midrule
0  & \QFourLowerZero & \QFourUpperZero & \QFourMidZero & \QFourMedianZero\\
6  & \QFourLowerSix & \QFourUpperSix & \QFourMidSix & \QFourMedianSix\\
12 & \QFourLowerTwelve & \QFourUpperTwelve & \QFourMidTwelve & \QFourMedianTwelve\\
24 & \QFourLowerTwentyFour & \QFourUpperTwentyFour & \QFourMidTwentyFour & \QFourMedianTwentyFour\\
36 & \QFourLowerThirtySix & \QFourUpperThirtySix & \QFourMidThirtySix & \QFourMedianThirtySix\\
\bottomrule
\end{array}
\]
On impact, the output-growth response is only partially pinned down:
\([\QFourLowerZero,\QFourUpperZero]\). This wide interval is exactly the point of
sign-restriction identification: many different rotations are observationally
equivalent and satisfy the same qualitative monetary-policy restrictions.

\subsection*{(ii)}
For the Bayesian thought experiment, we treat the reduced-form VAR estimates as fixed and
interpret Haar measure, truncated to the sign-restricted set, as the posterior over \(Q\).
The posterior median and equal-tailed 68\% credible set are therefore the 50th, 16th, and
84th percentiles of the retained \(\theta_\ell(q)\) draws. At the same selected horizons:
\[
\begin{array}{lrrr}
\toprule
\text{Horizon} & \text{68\% lower} & \text{Median} & \text{68\% upper}\\
\midrule
0  & \QFourCredLowerZero & \QFourMedianZero & \QFourCredUpperZero\\
6  & \QFourCredLowerSix & \QFourMedianSix & \QFourCredUpperSix\\
12 & \QFourCredLowerTwelve & \QFourMedianTwelve & \QFourCredUpperTwelve\\
24 & \QFourCredLowerTwentyFour & \QFourMedianTwentyFour & \QFourCredUpperTwentyFour\\
36 & \QFourCredLowerThirtySix & \QFourMedianThirtySix & \QFourCredUpperThirtySix\\
\bottomrule
\end{array}
\]
These credible sets summarize the arbitrary Haar weighting over rotations. They are not
frequentist confidence sets for the identified set and, as we know, they
need not cover the full range of sign-restriction-consistent responses.

\subsection*{(iii)}
Figure~\ref{fig:q4-sign-restricted-irf} plots the identified-set bounds, posterior median,
68\% credible set, and identified-set midpoint. The posterior median is not close to the
center of the identified set at every horizon. The largest absolute gap between the median
and midpoint is \(\QFourMaxMedianGap\) annualized percentage points, occurring at horizon
\(\QFourMaxMedianGapHorizon\). For example, at horizon 12 the identified-set midpoint is
\(\QFourMidTwelve\), while the posterior median is \(\QFourMedianTwelve\). This illustrates
that the truncated Haar prior supplies a weighting inside the identified
set rather than a neutral estimate of its center.

\begin{figure}[htbp]
\centering
\includegraphics[width=0.82\textwidth]{\QFourFigurePath}
\caption{Output-growth response to a one-standard-deviation contractionary monetary
policy shock under the sign restrictions that inflation falls and the real rate rises at
horizons 0, 1, and 2.}
\label{fig:q4-sign-restricted-irf}
\end{figure}

\end{document}
